% Part: sets-functions-relations
% Chapter: sets
% Section: proofs-about-sets

\documentclass[../../../include/open-logic-section]{subfiles}

\begin{document}

\olfileid[hi]{sfr}{set}{prf}
\olsection{समुच्चयों के बारे में प्रमाण}

\begin{editorial}
इस खंड का स्थान \verb|content/method/proofs| ने ले लिया है और इसे डिफ़ॉल्ट
रूप से इस अध्याय से हटा दिया गया है।
\end{editorial}

\begin{explain}
अब तक प्रस्तुत समुच्चय और संकेतन हमें समुच्चय निर्दिष्ट करने तथा उनके बीच
संबंध व्यक्त करने के सुविधाजनक संक्षिप्त रूप देते हैं। प्रायः ऐसे संबंधों के
बारे में दावे सिद्ध करना भी आवश्यक होगा। यदि आप गणितीय प्रमाणों से परिचित
नहीं हैं, तो यह आपके लिए नया हो सकता है। अतः हम एक सरल उदाहरण को
चरण-दर-चरण देखेंगे। हम सिद्ध करेंगे कि किन्हीं समुच्चयों $X$ और $Y$ के लिए
सदा $X \cap (X \cup Y) = X$ होता है। समुच्चयों के बीच ऐसी समानता कैसे सिद्ध
करते हैं? याद करें कि समुच्चय केवल अपने !!{element}s से निर्धारित होते हैं,
अर्थात् समुच्चय तभी और केवल तभी समान हैं जब उनके !!{element} समान हों। अतः
इस मामले में हमें सिद्ध करना है कि (a) $X \cap (X \cup Y)$ का प्रत्येक
!!{element}, $X$ का भी !!{element} है और इसके विपरीत, (b) $X$ का प्रत्येक
!!{element}, $X \cap (X \cup Y)$ का भी !!{element} है। दूसरे शब्दों में हम
दोनों दिखाते हैं: (a) $X \cap (X \cup Y) \subseteq X$ और (b) $X \subseteq
X \cap (X \cup Y)$।

``$X \cap (X \cup Y)$ का प्रत्येक !!{element}~$z$, $X$ का भी !!{element}
है'' जैसे सामान्य दावे को पहले कोई मनमाना $z \in X \cap (X \cup Y)$ दिया
हुआ मानकर और इस मान्यता से $z \in X$ सिद्ध करके सिद्ध किया जाता है। आप इस
प्रतिरूप को ``सामान्य सशर्त प्रमाण'' के नाम से जानते होंगे। इस प्रमाण में हमें
मान्यता और निष्कर्ष में आई परिभाषाओं का उपयोग भी करना होगा, उदाहरणतः इस मामले
में ``$\cap$'' और ``$\cup$'' की। अतः स्थिति (a) का तर्क इस प्रकार होगा:

\begin{quote}
(a) पहले हम दिखाना चाहते हैं कि $X \cap (X \cup Y) \subseteq X$, अर्थात्
$\subseteq$ की परिभाषा से, किसी भी~$z$ के लिए यदि $z \in X \cap (X \cup
Y)$, तो $z \in X$। अतः मान लें कि $z \in X \cap (X \cup Y)$। चूँकि $z$
दो समुच्चयों के सर्वनिष्ठ का !!{element} तभी और केवल तभी है जब वह दोनों
समुच्चयों का !!{element} हो, हम निष्कर्षित कर सकते हैं कि $z \in X$ और $z
\in X \cup Y$ भी। विशेषतः $z \in X$, और हमें यही दिखाना था।
\end{quote}

इससे प्रमाण का पहला आधा पूरा होता है। ध्यान दें कि अंतिम चरण में हमने यह
तथ्य उपयोग किया कि यदि किसी मान्यता से कोई संयोजन ($z \in X$ और $z \in X
\cup Y$) निष्पन्न होता है, तो उसी मान्यता से प्रत्येक संयोज्य निष्पन्न होता
है। आप इस नियम को ``संयोजन विलोपन'', या $\land$Elim, के नाम से जानते होंगे।
अब (b) सिद्ध करें:

\begin{quote}
(b) अब हम सिद्ध करते हैं कि $X \subseteq X \cap (X \cup Y)$, अर्थात्
$\subseteq$ की परिभाषा से, किसी भी~$z$ के लिए यदि $z \in X$, तो $z \in X
\cap (X \cup Y)$ भी। मान लें $z \in X$। यह दिखाने के लिए कि $z \in X
\cap (X \cup Y)$, हमें (``$\cap$'' की परिभाषा से) दिखाना है कि (i) $z \in
X$ और (ii) $z \in X \cup Y$ भी। यहाँ (i) केवल हमारी मान्यता है, इसलिए आगे
कुछ सिद्ध नहीं करना। (ii) के लिए याद करें कि $z$ समुच्चयों के किसी सम्मिलन
का !!{element} तभी और केवल तभी है जब वह उन समुच्चयों में कम-से-कम एक का
!!{element} हो। चूँकि $z \in X$ और $X \cup Y$, $X$ तथा $Y$ का सम्मिलन है,
यहाँ ऐसा ही है। अतः $z \in X \cup Y$। हमने (i) $z \in X$ और (ii) $z \in
X \cup Y$ दोनों दिखाए हैं; फलतः ``$\cap$'' की परिभाषा से $z \in X \cap (X
\cup Y)$।
\end{quote}

यह कुछ लंबा था, किंतु इससे स्पष्ट होता है कि हम समुच्चयों और उनके संबंधों के
बारे में कैसे तर्क करते हैं। सामान्यतः हम इतना स्पष्ट विवरण नहीं देते;
विशेषतः सभी परिभाषाएँ दोहराना आवश्यक नहीं होता। किसी अधिक उन्नत पाठ में हमारे
परिणाम का प्रमाण कहीं अधिक संक्षिप्त होगा। वह कुछ ऐसा दिख सकता है।
\end{explain}


\begin{prop}[अवशोषण]
सभी समुच्चयों $X$, $Y$ के लिए,
\[
X \cap (X \cup Y) = X
\]
\end{prop}

\begin{proof}
(a) मान लें $z \in X \cap (X \cup Y)$। तब $z \in X$, अतः $X \cap (X
  \cup Y) \subseteq X$।

(b) अब मान लें $z \in X$। तब $z \in X \cup Y$ भी और इसलिए $z \in X
\cap (X \cup Y)$ भी। अतः $X \subseteq X \cap (X \cup Y)$।
\end{proof}

\begin{prob}
विस्तार से सिद्ध कीजिए कि $X \cup (X \cap Y) = X$। फिर छोटा, संक्षिप्त
प्रमाण दीजिए। (संकेत: $X \cup (X \cap Y) \subseteq X$ दिशा के लिए आपको
स्थितियों द्वारा प्रमाण, अर्थात् $\lor$Elim, चाहिए होगा।)
\end{prob}

\end{document}
